\unitopen{33}
\begin{candos}
\item napisać list motywacyjny i e-mail aplikacyjny po niemiecku;
\item odpowiedzieć na dziesięć typowych pytań z rozmowy kwalifikacyjnej;
\item opisać mocne i słabe strony bez klisz;
\item podać oczekiwania finansowe i zapytać o warunki.
\end{candos}

\section{Słownictwo}
\vocabtable{33}

\section{List motywacyjny (\de{Anschreiben})}
Struktura na jedną stronę: (1) stanowisko i źródło ogłoszenia, (2) Twoje kwalifikacje związane z wymaganiami, (3) dlaczego ta firma,
(4) dostępność i zakończenie.

\begin{center}\small
\begin{tabular}{@{}p{0.2\linewidth}p{0.72\linewidth}@{}}
\toprule
\textsf{\textbf{Część}} & \textsf{\textbf{Przykład}}\\
\midrule
Wprowadzenie & \de{Mit großem Interesse habe ich Ihre Stellenanzeige als Senior Backend Engineer gelesen.}\\
Kwalifikacje & \de{In den letzten sechs Jahren habe ich Web-APIs mit .NET und Azure entwickelt und ein Team von drei Entwicklern geleitet.}\\
Motywacja & \de{Ihr Unternehmen reizt mich, weil es komplexe Systeme mit hohem Qualitätsanspruch baut.}\\
Zakończenie & \de{Ich freue mich auf ein persönliches Gespräch. Ich stehe ab dem 1.\ März zur Verfügung.}\\
\bottomrule
\end{tabular}
\end{center}

\section{Rozmowa kwalifikacyjna}
Najczęstsze pytania: \de{Erzählen Sie etwas über sich.} \de{Warum möchten Sie bei uns arbeiten?} \de{Was sind Ihre Stärken
und Schwächen?} \de{Beschreiben Sie ein Projekt, auf das Sie stolz sind.} \de{Wie gehen Sie mit Konflikten um?} \de{Wo sehen Sie sich in
fünf Jahren?} \de{Wie sind Ihre Gehaltsvorstellungen?}

\begin{zasada}[Odpowiedź w trzech krokach: sytuacja, działanie, wynik]
\de{In meinem letzten Projekt hatten wir eine hohe Latenz. Ich habe die Abfragen analysiert und einen Cache eingeführt.
Dadurch ist die Antwortzeit um die Hälfte gesunken.} Mów w Perfekcie o zakończonych działaniach i konkretnie, nie ogólnikowo.
\end{zasada}

\begin{falle}
Polska kultura rozmowy bywa skromna, a szwajcarska wymaga konkretów i spokojnej pewności siebie. Błąd: \de{Ich bin sehr gut
in allem} (ogólnik). Lepiej: \de{Ich bin stark in der Fehleranalyse; dafür habe ich im letzten Jahr eine Monitoring-Lösung gebaut.}
Przy słabej stronie nie mów \de{Ich bin Perfektionist} (klisza), tylko prawdziwą słabość i środek zaradczy.
\end{falle}

\begin{szwajcaria}
Szwajcarskie oferty podają często przedział płacy w CHF brutto rocznie. Rynek IT w Zurychu wypłaca wysokie wynagrodzenia, ale
koszty życia są wysokie, więc porównuj netto i koszty. Pytanie \de{Was sind Ihre Lohnvorstellungen?} traktuj jak
standard: przygotuj zakres, a nie jedną liczbę. Standardowe dodatki: \de{13.\ Monatslohn}, \de{Pensionskasse}, \de{Ferien}.
Obywatel UE potrzebuje na pracę zwykle pozwolenia \de{Aufenthaltsbewilligung B} \verify{(zasady dla obywateli UE/EFTA wynikają z umowy
o swobodnym przepływie osób, sprawdź aktualne wymogi w SEM i w kantonie)}.
\end{szwajcaria}

\begin{praca}
Mail aplikacyjny: \de{Sehr geehrte Frau Keller, ich bewerbe mich um die Stelle als Senior Backend Engineer. Im Anhang finden Sie meinen
Lebenslauf und mein Motivationsschreiben. Ich freue mich auf Ihre Rückmeldung. Freundliche Grüsse, Piotr Nowak.}
\end{praca}

\section{Dialogi}
\textbf{Dialog 1 --- rozmowa kwalifikacyjna (fragment)}
\begin{dialog}
\turn{A}{Warum möchten Sie bei uns arbeiten?}{Dlaczego chce Pan u nas pracować?}
\turn{B}{Weil Ihr Unternehmen komplexe Systeme baut und großen Wert auf Qualität legt. Ich möchte in einem Team arbeiten, von dem ich lernen kann.}{Ponieważ Państwa firma buduje złożone systemy i przywiązuje dużą wagę do jakości. Chcę pracować w zespole, od którego mogę się uczyć.}
\turn{A}{Was sind Ihre Stärken?}{Jakie są Pana mocne strony?}
\turn{B}{Ich analysiere Probleme gründlich und arbeite gern im Team. Eine Schwäche ist, dass ich manchmal zu lange an Details arbeite.}{Analizuję problemy dokładnie i lubię pracować w zespole. Słabością jest, że czasem za długo pracuję nad szczegółami.}
\turn{A}{Wie sind Ihre Gehaltsvorstellungen?}{Jakie są Pana oczekiwania finansowe?}
\turn{B}{Ich stelle mir ein Gehalt zwischen 130.000 und 150.000 Franken vor. Das hängt auch vom Gesamtpaket ab.}{Wyobrażam sobie wynagrodzenie między 130 000 a 150 000 franków. Zależy to też od całego pakietu.}
\end{dialog}

\section{Ćwiczenia}
\begin{cwiczenia}
\exitem{Napisz pierwsze zdanie listu na stanowisko Backend Engineer (\de{Mit großem Interesse \dots}).}{\de{Mit großem Interesse habe ich Ihre Stellenanzeige als Backend Engineer gelesen.}}
\exitem{Powiedz: „Aplikuję na stanowisko programisty.”}{\de{Ich bewerbe mich um die Stelle als Entwickler.}}
\exitem{Odpowiedz: \de{Warum möchten Sie bei uns arbeiten?} (dwa zdania z \de{weil}).}{np.\ \de{Weil Ihr Unternehmen interessante Projekte hat und weil ich im Team lernen möchte.}}
\exitem{Podaj jedną mocną stronę z przykładem (Perfekt).}{np.\ \de{Ich bin gut in der Fehleranalyse. Im letzten Projekt habe ich einen kritischen Bug gefunden.}}
\exitem{Podaj jedną słabość i środek zaradczy.}{np.\ \de{Ich verliere mich manchmal in Details. Deshalb setze ich mir Zeitlimits.}}
\exitem{Podaj oczekiwania finansowe zakresem.}{np.\ \de{Ich stelle mir ein Jahresgehalt zwischen X und Y Franken vor.}}
\exitem{Zapytaj o warunki pracy: wymiar i urlop.}{\de{Wie viele Urlaubstage gibt es und wie sieht die Arbeitszeit aus?}}
\exitem{Popraw: \de{Ich bewerbe mich für die Stelle auf.}}{\de{Ich bewerbe mich um die Stelle.}}
\exitem{Przetłumacz: „Czy mogę zadać pytanie o zespół?”}{\de{Darf ich eine Frage zum Team stellen?}}
\exitem{Przetłumacz: „Chętnie przyjdę na rozmowę.”}{\de{Ich komme gern zu einem Gespräch.}}
\end{cwiczenia}

\section{Trening słownictwa (generowany)}
\drill{33}

\weekend{Napisz pełny list motywacyjny (jedna strona) na wymarzone stanowisko w Zurychu i powiedz na głos odpowiedzi na dziesięć
pytań z rozmowy. Nagraj 5 minut i poproś o korektę form i szyku.}
